\documentclass[12pt,a4paper,oneside]{report}
\usepackage[T1]{fontenc}
\usepackage[utf8]{inputenc}
\usepackage[french]{babel}
\usepackage{lmodern,microtype}
\usepackage[a4paper,margin=2.35cm]{geometry}
\usepackage{setspace}
\onehalfspacing
\setlength{\parindent}{0pt}
\setlength{\parskip}{0.5\baselineskip}
\setlength{\emergencystretch}{3em}
\clubpenalty=10000
\widowpenalty=10000
\interlinepenalty=10000
\raggedbottom
\usepackage{amsmath,amssymb,booktabs,longtable,array}
\usepackage{xcolor,listings,hyperref}
\hypersetup{colorlinks=true,linkcolor=blue!50!black,urlcolor=blue!50!black,
pdftitle={Cours 13 -- Coder un reseau neuronal en partant de zero},pdfauthor={Projet IPR}}
\lstset{language=Python,basicstyle=\ttfamily\footnotesize,breaklines=true,
columns=fullflexible,keepspaces=true,showstringspaces=false,
keywordstyle=\color{blue!55!black},commentstyle=\color{green!35!black},
frame=single,aboveskip=1em,belowskip=1em}
\newcommand{\retenir}[1]{\paragraph{À retenir.} #1}
\title{Cours 13 -- Coder un réseau neuronal\\\large Des premières équations à un prototype complet de restauration d'image}
\author{Projet IPR -- Documentation Scholar}
\date{31 octobre 2026\\\small Aucun prérequis en apprentissage automatique. Exemples originaux exécutables.}

\begin{document}
\maketitle
\tableofcontents

\chapter{Ce que nous allons construire}

Un réseau neuronal est un programme paramétrable. Ses paramètres numériques,
appelés poids et biais, sont ajustés à partir d'exemples. Il ne comprend pas
automatiquement la photographie et ne découvre pas une vérité physique sans
données appropriées. Il apprend des relations qui peuvent être utiles, mais
aussi des biais, des défauts et des raccourcis présents dans les exemples.

Notre objectif est de comprendre chaque étape : représenter les données,
calculer une prédiction, mesurer l'erreur, calculer ses dérivées, modifier les
paramètres, puis vérifier que le résultat fonctionne sur des exemples nouveaux.
Nous commencerons avec quelques nombres, puis passerons aux images et à un U-Net.

\section{Trois significations de « partir de zéro »}

On peut écrire les calculs sans bibliothèque, programmer une architecture connue
avec PyTorch, ou entraîner ses propres poids sans modèle préentraîné. Ces trois
choix sont différents. Nous ferons les deux premiers, puis un petit ajustement
synthétique pour comprendre le troisième. Un entraînement photographique sérieux
reste un projet de données et de qualification distinct.

Une architecture éprouvée n'impose pas de reprendre les poids de ses auteurs.
Inversement, posséder ses poids n'assure ni leur qualité ni des droits suffisants
sur les images ayant servi à les apprendre. Séparer architecture, implémentation,
poids, données et moteur d'exécution évite beaucoup de confusions.

\section{Vocabulaire initial}

\begin{longtable}{p{.25\textwidth}p{.67\textwidth}}
\toprule Terme & Sens concret \\ \midrule
Entrée & Nombres fournis au programme, par exemple une image bruitée.\\
Cible & Résultat souhaité pour un exemple, par exemple une référence moins bruitée.\\
Prédiction & Résultat actuellement calculé par le réseau.\\
Paramètre & Nombre modifiable pendant l'apprentissage : poids ou biais.\\
Hyperparamètre & Choix de conception : largeur, taux d'apprentissage, taille des crops.\\
Perte & Nombre qui mesure une erreur selon une règle choisie.\\
Gradient & Dérivées de la perte par rapport aux paramètres.\\
Entraînement & Calcul répété des prédictions, gradients et mises à jour.\\
Inférence & Utilisation des poids figés pour calculer une sortie.\\
Généralisation & Capacité à fonctionner sur des exemples non utilisés pour les réglages.\\
\bottomrule
\end{longtable}

\chapter{Préparer un atelier Python isolé}

Une variable nomme une valeur, une liste regroupe des valeurs, une boucle répète
une opération et une fonction rassemble un calcul réutilisable. Python permet
d'expérimenter ces notions sans compiler un exécutable. Les fichiers du cours
restent indépendants du moteur de production d'IPR.

\begin{lstlisting}
def prediction(x, weight, bias):
    return weight * x + bias

values = [-1.0, 0.0, 1.0]
for x in values:
    print(prediction(x, 2.0, 1.0))
\end{lstlisting}

\section{Environnement virtuel}

Un environnement virtuel contient ses propres paquets Python. Il évite de
modifier les dépendances d'un autre agent ou celles des campagnes alpha.
Depuis un dossier de travail dédié, les commandes suivantes illustrent une
installation CPU ; adapter le nom de l'interpréteur installé sur la machine.

\begin{lstlisting}[language={}]
python -m venv .venv-cours
.venv-cours\Scripts\python.exe -m pip install numpy==2.3.5
.venv-cours\Scripts\python.exe -m pip install torch==2.9.1+cpu --index-url https://download.pytorch.org/whl/cpu
\end{lstlisting}

Sous Linux/macOS, l'interpréteur de cet environnement se trouve généralement
dans \texttt{.venv-cours/bin/python}. Les versions ci-dessus correspondent au
témoin CPU de ce cours sur Windows/Python 3.14, pas à une qualification universelle.
Pour un GPU, choisir explicitement un paquet compatible avec le pilote et le
système : avoir une carte NVIDIA ne suffit pas à rendre CUDA utilisable.

\section{Exemples fournis}

Depuis la racine du worktree documentaire :
\begin{lstlisting}[language={}]
python docs/scholar/references/reseaux-exemples.py
python docs/scholar/references/reseaux-exemples.py --torch
\end{lstlisting}

La première commande utilise seulement la bibliothèque standard. La seconde
nécessite PyTorch dans l'interpréteur choisi. Aucun exemple ne télécharge de
poids, ne lit de photo, ne crée de fichier ou ne modifie un RAW. Le petit exemple
PyTorch est limité à deux threads CPU et à une image artificielle.

\chapter{Les mathématiques indispensables, sans raccourci}

\section{Scalaires, vecteurs, matrices et tenseurs}

Un scalaire est un nombre, comme une température. Un vecteur est une suite de
nombres : un pixel RGB peut être noté $x=(R,G,B)$. Une matrice est un tableau à
deux axes. Un tenseur généralise ces tableaux à plusieurs axes. Le mot ne signifie
pas à lui seul qu'il s'agit d'un réseau ou d'une opération compliquée.

Une image RGB de hauteur $H$ et largeur $W$ contient $3HW$ valeurs. En mémoire
intercalée, l'ordre peut être $H\times W\times3$. Les couches de convolution
PyTorch utilisent couramment $N\times C\times H\times W$, où $N$ est le nombre
d'images traitées ensemble et $C$ le nombre de canaux. Changer l'ordre des axes
demande une permutation, pas un simple changement de dimensions.

Le produit scalaire $w^Tx=\sum_iw_ix_i$ combine plusieurs entrées. Une couche
dense calcule $z=Wx+b$. Si $x$ a $d$ composantes et $z$ en a $m$, $W$ a
$m\times d$ coefficients et $b$ en a $m$. Il y a donc $md+m$ paramètres.

\section{Fonction, pente et dérivée}

Une fonction associe une sortie à une entrée. La dérivée décrit comment la
sortie change pour une petite variation de l'entrée. Pour $f(w)=w^2$,
$f'(w)=2w$. Une fonction avec plusieurs paramètres possède une dérivée par
paramètre : leur ensemble est le gradient.

La différence finie permet de vérifier une dérivée programmée :
\[
 f'(w)\approx\frac{f(w+\varepsilon)-f(w-\varepsilon)}{2\varepsilon}.
\]
Un $\varepsilon$ trop grand introduit une erreur d'approximation ; trop petit,
il expose aux arrondis. Ce contrôle sert à déboguer un petit calcul, pas à
entraîner efficacement des millions de poids.

\section{Moyenne, variance et probabilité}

La moyenne estime un niveau central. La variance mesure une dispersion :
$\operatorname{Var}(X)=\mathbb{E}[(X-\mathbb{E}X)^2]$. L'écart-type est sa
racine carrée. Une loi gaussienne est décrite par une moyenne et une variance,
mais le bruit d'un capteur réel n'est pas toujours gaussien, indépendant ou
uniforme sur toute l'image.

Une probabilité prédit une fréquence ou une incertitude sous des hypothèses.
Un débruiteur ne doit pas être confondu avec une mesure certaine de la valeur
réelle d'un pixel. Plusieurs images propres peuvent expliquer une observation
très bruitée ; aucun programme ne recrée une information détruite avec garantie.

\chapter{Premier apprentissage : retrouver une droite}

Supposons que des exemples suivent $y=2x+1$, mais que le programme ne connaisse
pas les coefficients. Nous choisissons $\hat y=wx+b$, initialisons $w=b=0$,
puis mesurons l'erreur quadratique moyenne :
\[
 L=\frac1n\sum_{i=1}^n(wx_i+b-y_i)^2.
\]
En posant $e_i=wx_i+b-y_i$, on obtient
\[
 \frac{\partial L}{\partial w}=\frac2n\sum_i e_ix_i,
 \qquad \frac{\partial L}{\partial b}=\frac2n\sum_i e_i.
\]
La descente de gradient effectue $w\leftarrow w-\eta\partial L/\partial w$ et
$b\leftarrow b-\eta\partial L/\partial b$. Le taux $\eta$ règle la longueur
du pas. Le signe moins cherche à descendre la perte, sans garantir de trouver
son meilleur minimum pour tout problème et tout taux.

\lstinputlisting[firstline=12,lastline=24]{references/reseaux-exemples.py}

Ce programme retrouve environ $w=2$ et $b=1$. Il contient déjà l'essentiel d'une
boucle d'apprentissage : données, prédiction, erreur, gradients, mise à jour.
Une architecture plus grande répète ces principes avec davantage de paramètres.

\section{Un pas calculé à la main}

Pour $x=2$, $y=5$, $w=2$, $b=2$, la prédiction vaut 6. Avec
$L=\tfrac12(\hat y-y)^2$, la perte vaut 0,5, les gradients valent 2 pour $w$
et 1 pour $b$. À $\eta=0,1$, les nouveaux paramètres sont 1,8 et 1,9. La prédiction
devient 5,5 et la perte 0,125. Le facteur $1/2$ explique l'absence de facteur 2
dans cette variante ; toujours lire la définition exacte de la perte.

\chapter{Du neurone au réseau multicouche}

Un neurone calcule d'abord $z=w^Tx+b$, puis une activation $a=\phi(z)$.
Sans activation non linéaire, plusieurs couches affines successives restent une
seule transformation affine : empiler des couches ne suffit donc pas.

\begin{longtable}{p{.22\textwidth}p{.32\textwidth}p{.37\textwidth}}
\toprule Activation & Formule & Utilité et limite \\ \midrule
ReLU & $\max(0,z)$ & Simple ; dérivée nulle côté négatif.\\
Sigmoïde & $1/(1+e^{-z})$ & Sortie entre 0 et 1 ; saturation aux extrêmes.\\
Tanh & $\tanh(z)$ & Sortie signée bornée ; peut saturer.\\
SiLU & $z/(1+e^{-z})$ & Activation lisse, utilisée dans le TP.\\
Identité & $z$ & Sortie sans borne imposée, utile au RGB signé.\\
\bottomrule
\end{longtable}

Une activation interne ReLU ne force pas nécessairement une sortie RGB positive :
la dernière couche peut combiner les caractéristiques avec des poids signés.
En revanche, une sigmoïde finale contraindrait réellement la sortie à $(0,1)$,
ce qui ne correspond pas au signal étendu d'IPR.

\section{Pourquoi XOR demande de la non-linéarité}

XOR vaut 1 si deux bits diffèrent, sinon 0. Une seule frontière droite ne sépare
pas ses quatre exemples. Une couche cachée peut construire plusieurs régions,
puis les combiner. Le TP entraîne huit neurones cachés avec tanh et une sortie
sigmoïde ; les gradients sont calculés à la main, sans PyTorch.

\lstinputlisting[firstline=38,lastline=82]{references/reseaux-exemples.py}

Ce petit problème n'est pas une photographie miniature : il sert à comprendre
les paramètres et la circulation des dérivées. Le nombre de neurones d'un
modèle image ne peut pas être choisi en extrapolant les quatre exemples XOR.

\chapter{Comprendre et programmer la rétropropagation}

\section{La règle de la chaîne}

Si $a=\phi(z)$ et $z=wx+b$, alors
$\partial L/\partial w=(\partial L/\partial a)\phi'(z)x$.
Chaque opération multiplie le gradient reçu par sa propre dérivée locale.
Un graphe de calcul garde les dépendances pour appliquer cette règle dans
l'ordre inverse du calcul direct.

Pour une couche dense, en notant $\delta=\partial L/\partial z$ :
\[
 \frac{\partial L}{\partial W}=\delta x^T,\qquad
 \frac{\partial L}{\partial b}=\delta,\qquad
 \frac{\partial L}{\partial x}=W^T\delta.
\]
Pour un lot, sommer ou moyenner les contributions selon la réduction de la perte.
Ne pas moyenner deux fois par erreur : cela change la taille des gradients.

\section{Un cas où les formules se simplifient}

La perte binaire est
$L=-[y\log(p)+(1-y)\log(1-p)]$. Avec $p=\operatorname{sigmoid}(z)$,
$\partial L/\partial z=p-y$. Pour tanh, la dérivée vaut $1-a^2$.
Ce sont les deux relations utilisées dans le TP XOR. Les produits utilisent
les poids du même état avant leur mise à jour ; les modifier au milieu de
la rétropropagation mélangerait deux versions du réseau.

\section{Ce que PyTorch automatise}

Un tenseur avec \texttt{requires\_grad=True} participe au suivi des gradients.
\texttt{loss.backward()} remplit les champs \texttt{.grad}. L'optimiseur les
utilise ensuite ; il ne calcule pas lui-même la vérité ou la fonction de perte.
Les gradients s'accumulent par défaut, d'où \texttt{zero\_grad()} à chaque pas,
sauf accumulation de lots explicitement voulue et normalisée.

\begin{lstlisting}
optimizer.zero_grad(set_to_none=True)
prediction = model(inputs)
loss = (prediction - target).square().mean()
loss.backward()
optimizer.step()
\end{lstlisting}

\texttt{model.eval()} change certains modules comme dropout ou BatchNorm,
mais ne désactive pas les gradients. \texttt{torch.no\_grad()} désactive leur
suivi, mais ne met pas automatiquement le modèle en mode évaluation. Pour une
inférence habituelle, employer les deux, avec des contrats d'entrée vérifiés.

\chapter{Choisir une perte et un optimiseur}

\section{L'erreur choisie définit ce que le réseau préfère}

La MSE pénalise fortement les grands écarts et conduit théoriquement à la
moyenne conditionnelle si le modèle et l'optimisation sont idéaux. Une ambiguïté
sur les détails peut donc produire du lissage. La MAE conduit à une médiane
conditionnelle ; elle n'assure pas davantage la préservation de chaque texture.

Une perte de Charbonnier moyenne $\sqrt{e^2+\varepsilon^2}$ et fournit un
compromis lisse. Des termes de gradients ou de couleur peuvent compléter la
perte pixel, mais leurs coefficients changent l'objectif : les qualifier par
ablation plutôt que les additionner parce qu'ils semblent utiles.

Pour un débruiteur scientifique, une perte perceptuelle ou adversariale ne doit
pas devenir un moyen silencieux d'inventer de la matière. Une image agréable
peut être moins fidèle. Choisir d'abord fidélité, biais photométrique et détails
connus, puis une revue visuelle sur de nouvelles scènes.

\section{SGD, momentum et Adam}

SGD applique une mise à jour à partir d'un lot. Le momentum conserve une moyenne
des gradients pour stabiliser la direction. Adam suit deux moyennes :
\[
 m_t=\beta_1m_{t-1}+(1-\beta_1)g_t,\quad
 v_t=\beta_2v_{t-1}+(1-\beta_2)g_t^2.
\]
Après corrections de biais,
$\theta_{t+1}=\theta_t-\eta\hat m_t/(\sqrt{\hat v_t}+\varepsilon)$.
Le carré et la division sont composante par composante. Adam n'empêche pas
le surapprentissage ; AdamW traite la décroissance des poids séparément de
l'adaptation du pas. Comparer un petit nombre de choix explicitement définis.

\section{Initialisation et stabilité}

Initialiser tous les neurones cachés de façon identique maintient des symétries
qui limitent leur apprentissage. Les initialisations aléatoires contrôlées,
comme Xavier ou He selon les activations, évitent ce piège. Fixer la graine
permet de répéter un essai, sans garantir une identité numérique entre GPU.

Le TP initialise seulement la dernière couche résiduelle à zéro : la sortie
initiale vaut exactement l'entrée, tandis que le corps est aléatoire. Au premier
pas, le gradient apprend surtout cette tête ; une fois ses poids non nuls, les
couches précédentes reçoivent leur gradient. Ce n'est pas l'initialisation de
tout le réseau à zéro.

\chapter{Construire les données avant de grossir le modèle}

Un dataset associe les entrées aux cibles et métadonnées. Un chargeur compose
les lots. Un crop est une région de l'image, pas une nouvelle scène indépendante.
Les transformations géométriques doivent être identiques pour l'entrée et sa
cible. Un décalage d'un pixel peut enseigner du flou au lieu du débruitage.

\section{Trois groupes et deux niveaux de décision}

Le jeu d'entraînement met à jour les poids. La validation aide à choisir les
hyperparamètres et le checkpoint. Le test final évalue le choix figé ; il ne
sert pas à régler une nouvelle variante. Si l'on consulte sans cesse le même
test pour choisir, il devient indirectement une validation.

Séparer par scènes, séances et, pour tester une généralisation inter-boîtiers,
par modèles de capteurs. Mettre deux images d'une rafale dans deux groupes ne
crée pas cette indépendance. Conserver identités, empreintes, provenance,
autorisations et raisons d'exclusion dans une mémoire versionnée.

\section{D'où vient une cible propre ?}

Une moyenne de prises alignées atténue un bruit temporel indépendant ; elle
conserve le bruit fixe commun et peut lisser le mouvement. Une photo basse ISO
en pose longue est un candidat, pas une vérité identique par définition :
éclairage, saturation, exposition et signal capteur doivent être raccordés.
Une image Lightroom débruitée est le résultat d'un autre algorithme, non une
mesure certaine de la scène.

Le bruit synthétique permet de connaître la cible, mais un bruit gaussien
ajouté après dématriçage ne reproduit pas tout un capteur. Mélanger données
réelles qualifiées et simulation documentée est une stratégie à étudier,
sans confondre leurs niveaux de preuve.

\chapter{Des nombres aux images : la convolution}

Une couche dense appliquée à tous les pixels aurait énormément de paramètres
et ne partagerait pas naturellement les motifs spatiaux. Une convolution utilise
un petit ensemble de poids qui glisse sur l'image. Le même détecteur local peut
ainsi répondre à un bord situé en plusieurs endroits.

Dans les bibliothèques, l'opération dite convolution est généralement une
corrélation croisée, sans retournement du noyau :
\[
 Y_{o,h,w}=b_o+\sum_{c,u,v}K_{o,c,u,v}X_{c,h+u,w+v}.
\]
Les conventions d'origine, padding et stride complètent cette formule. Avec
un noyau $k\times k$, $C_{in}$ canaux et $C_{out}$ sorties, le nombre de
paramètres est $k^2C_{in}C_{out}+C_{out}$ si le biais est activé.

Pour $3$ entrées, $8$ sorties et $k=3$, on obtient 224 paramètres. Pour les six
canaux du TP, la première couche $6\to8$ en compte 440. Une carte de
caractéristiques n'est plus un canal rouge ou bleu : c'est une représentation
apprise, dont la signification n'est pas imposée individuellement.

\section{Dimensions et contexte}

Avec padding $p$, dilatation $d$ et stride $s$, une dimension devient
\[
 H_{out}=\left\lfloor\frac{H+2p-d(k-1)-1}{s}+1\right\rfloor.
\]
Deux convolutions $3\times3$, stride 1, voient un voisinage théorique
$5\times5$. Des réductions de résolution augmentent plus vite le contexte,
mais peuvent perdre des détails. Le contexte théorique ne signifie pas que
tous les pixels ont la même influence effective.

Le padding nul invente des valeurs en bordure ; réflexion et réplication ont
leurs propres hypothèses. Le tuilage doit conserver un halo cohérent. Dès que
le réseau emploie des statistiques globales, un halo fini ne garantit plus la
même sortie qu'une image entière.

\chapter{Concevoir un petit U-Net complet}

U-Net associe un encodeur qui réduit la résolution à un décodeur qui la remonte.
Les connexions latérales apportent des caractéristiques fines de l'encodeur
au décodeur. Le principe vient d'un travail de segmentation ; son adaptation
à la restauration demande un objectif et une sortie appropriés.

\begin{center}
\begin{tabular}{p{.20\textwidth}p{.37\textwidth}p{.33\textwidth}}
\toprule Étape & Dimensions pour une entrée $32\times32$ & Rôle \\ \midrule
Entrée & $1\times6\times32\times32$ & RGB et sigma par canal\\
Encodeur1 & $1\times8\times32\times32$ & Caractéristiques fines\\
Encodeur2 & $1\times16\times16\times16$ & Contexte intermédiaire\\
Centre & $1\times32\times8\times8$ & Contexte réduit\\
Décodeur2 & $1\times16\times16\times16$ & Fusion avec encodeur2\\
Décodeur1 & $1\times8\times32\times32$ & Fusion avec encodeur1\\
Tête & $1\times3\times32\times32$ & Résidu signé\\
\bottomrule
\end{tabular}
\end{center}

La concaténation additionne les nombres de canaux :32+16=48 au premier
décodeur, puis16+8=24 au second. Elle n'est pas une addition pixel à pixel.
Le rééchantillonnage remonte les dimensions mais ne recrée pas à lui seul
l'information perdue ; les convolutions et connexions latérales complètent
le calcul.

\section{Pourquoi prédire un résidu}

On peut prédire l'image propre directement ou une correction $r_\theta$ :
$\hat x=y-r_\theta(y,\sigma)$. Le résidu n'est pas automatiquement une
réalisation exacte du bruit physique. La formulation facilite une identité
initiale et une interface de force, mais ne garantit ni absence de biais ni
protection de chaque détail.

La largeur 8 et les deux réductions du TP donnent 29\,995 paramètres.
C'est un petit témoin pédagogique et technique, pas une taille minimale
démontrée pour les photos Canon. Il ne faut pas confondre le nombre de couches,
le nombre de canaux, le nombre de poids et le nombre de pixels de contexte.

\chapter{Coder et entraîner ce U-Net avec PyTorch}

Voici le laboratoire complet : imports locaux, définition des blocs, classe,
données synthétiques, optimiseur, calcul direct et rétropropagation. La fonction
peut être exécutée avec l'option \texttt{--torch} du script fourni.

\lstinputlisting[firstline=85,lastline=142]{references/reseaux-exemples.py}

\section{Lire le code sans magie}

Une sous-classe de \texttt{nn.Module} enregistre ses couches et leurs paramètres.
\texttt{forward} décrit le calcul. Appeler \texttt{model(image,sigma)} passe
par ce calcul et les mécanismes du framework. \texttt{model.parameters()}
fournit les poids à l'optimiseur ; les tenseurs de l'image ne deviennent pas
des paramètres parce qu'ils traversent le réseau.

\texttt{torch.cat} réunit RGB et cartes sigma. Le pooling réduit la taille.
\texttt{interpolate} rejoint les dimensions exactes des connexions latérales,
y compris lorsque certaines dimensions sont impaires. Les couches finales
ne comportent ni sigmoïde ni clamp. Le modèle travaille en flottants signés.

\section{Ce que ce laboratoire prouve et ne prouve pas}

Une perte qui baisse sur une seule rampe prouve au mieux que des gradients et
une mise à jour fonctionnent dans ce cas. Elle ne montre pas de généralisation,
de fidélité des cheveux ou de performance sur une nuit lunaire. Un modèle
peut mémoriser une image et échouer sur toutes les autres.

Pour une vraie boucle, remplacer l'image fixe par des lots du dataset, garder
un jeu de validation indépendant, calculer les métriques sous
\texttt{no\_grad}, et sélectionner le checkpoint par la règle fixée avant
le test final. Les augmentations, normalisations et masques font partie du
modèle au sens opérationnel : les conserver avec les poids.

\chapter{Cas IPR : couleur, bruit et valeurs négatives}

Un réseau ne sait pas spontanément si un triplet représente du sRGB encodé,
du Rec.2020 linéaire ou des axes caméra. Le même nombre n'a pas le même sens
dans ces domaines. Définir ordre RGB/BGR, gains, transfert, blanc, noir, unités,
orientation et stade du pipeline est aussi important que choisir les couches.

La cible étudiée pour IPR est un RGB caméra déjà balancé, après RCD/médiane,
avant optique et DCP lorsque la source pré-optique est disponible. Les valeurs
négatives sous le noir sont des observations utiles, pas nécessairement des
erreurs de fichier. Des valeurs au-dessus de 1 peuvent également être finies
et recevables dans un domaine étendu.

\section{Modèle physique et conditionnement}

Dans une approximation capteur, pour un signal attendu $s\geq0$,
$\operatorname{Var}(n\mid s)\approx as+b$. Le terme de signal provient
notamment du comptage des photons ; le terme constant représente une partie
du bruit de lecture. Après gains, matrices et dématriçage, corrélations et
variance changent : trois cartes sigma ne décrivent pas toute la covariance.

Donner une carte de bruit au réseau peut rendre ses décisions plus explicites,
mais cette carte doit être calibrée ou son incertitude représentée. L'ISO seul
n'est pas une puissance de bruit universelle. Changer de boîtier peut changer
bruit fixe, réponse spectrale, niveaux et traitements intégrés.

\section{Ce que nous ne déléguerons pas silencieusement}

Les conversions couleur connues et les profils restent des transformations
déterministes. Le réseau propose une correction, pas un nouveau choix créatif
de couleur. Mesurer séparément biais dans les noirs, conservation des gris,
écarts des couleurs saturées et détail. L'absence de clamp final autorise
les valeurs signées, mais ne suffit pas à apprendre leur bon traitement.

La force zéro contourne exactement le modèle. Une interpolation
$x_t=y+t(\hat x-y)$ donne une commande simple, pas une garantie de qualité
monotone ; tester les valeurs intermédiaires. Ne pas débruiter cumulativement
la sortie précédente à chaque mouvement du curseur : repartir de la source
immuable et de la recette courante.

\chapter{Passer du laboratoire à un projet complet}

\begin{longtable}{p{.25\textwidth}p{.67\textwidth}}
\toprule Partie & Responsabilité \\ \midrule
Contrat & Axes, unités, domaine couleur, valeurs valides et erreurs.\\
Données & Identités, partitions, références, masques et droits.\\
Modèle & Architecture, configuration, initialisation et poids.\\
Entraînement & Lots, perte, optimiseur, état de reprise et graines.\\
Évaluation & Métriques, scènes indépendantes, contrôles et revue visuelle.\\
Inférence & Conversion des buffers, tuilage, ressources et annulation.\\
Déploiement & Runtime, licences, empreintes et compatibilité des opérateurs.\\
Suivi & Versions, résultats JSON, limites et décisions de validation.\\
\bottomrule
\end{longtable}

Un checkpoint de reprise contient généralement poids, état de l'optimiseur,
compteur d'étapes, configuration et états aléatoires. Une livraison d'inférence
n'a pas besoin de tout cet état. Ne pas exécuter un pickle ou un script inconnu
pour obtenir des poids ; provenance, chargement limité aux tenseurs et
empreintes doivent être contrôlés.

Exporter vers ONNX ne garantit pas une identité numérique ou la prise en
charge de chaque opérateur. Comparer les sorties de référence et exportées
sur les mêmes entrées, y compris tailles impaires, négatifs, constantes et
bords. Le passage FP32 vers FP16 ou une quantification int8 demande une
qualification propre ; des économies mémoire ne justifient pas une dérive cachée.

\section{RAM, VRAM et temps}

Les poids ne sont qu'une partie de la mémoire. Ajouter activations, gradients,
états de l'optimiseur, lots et buffers de transfert. Adam conserve notamment
des états par paramètre. L'entraînement coûte donc plus de mémoire que
l'inférence du même réseau. Une image RGB float32 de 32 Mpx occupe 384 Mo
décimaux avant le moindre réseau.

Commencer avec de petits crops et un seul lot, mesurer les pics, puis augmenter
progressivement. Le disque peut conserver les RAW et un cache borné, mais
ne remplace pas directement la VRAM d'une opération. Fixer un budget commun
aux tâches pour éviter plusieurs réseaux qui pensent chacun disposer de toute
la machine.

\chapter{Tester, diagnostiquer et décider}

\section{Ordre conseillé}

Vérifier d'abord dimensions et unités. Faire ensuite fonctionner un petit lot,
contrôler les gradients puis surajuster volontairement quelques exemples pour
détecter une erreur de raccordement. Passer enfin à des données indépendantes
avec un protocole figé. Un échec dès les petites étapes évite un entraînement
coûteux dont les résultats seraient ininterprétables.

\begin{longtable}{p{.30\textwidth}p{.62\textwidth}}
\toprule Symptôme & Hypothèses à vérifier \\ \midrule
Perte NaN & Valeur invalide, division/logarithme, pas trop grand, débordement.\\
Perte constante & Gradient coupé, poids non donnés à l'optimiseur, tête non reliée.\\
Belle image très lisse & Cible floue, perte favorisant la moyenne, bruit mal simulé.\\
Couleur modifiée & RGB/BGR, transfert, gains, normalisation ou cible incohérents.\\
Bon entraînement, mauvais test & Fuite de données, surapprentissage, domaine non couvert.\\
Coutures & Padding, halo, fusion ou statistiques spatiales du réseau.\\
RAM croissante & Graphes conservés, sorties accumulées, chargeurs ou caches non bornés.\\
\bottomrule
\end{longtable}

\section{Mesures et validation}

Avec une référence admissible, MSE et PSNR sont utiles ;
$\operatorname{PSNR}=10\log_{10}(M^2/\operatorname{MSE})$, où $M$ doit être
défini. Pour un domaine signé étendu, annoncer $M=1$ par convention ne signifie
pas que toutes les valeurs sont bornées à1. SSIM et mesures perceptuelles
ajoutent des indications, mais aucune moyenne unique ne couvre la fidélité.

Sans référence propre, une baisse de l'énergie haute fréquence peut être
du lissage de texture. Employer des masques fixes, des détails connus,
des comparaisons visuelles à taille native et des mesures de biais/couleur.
Une rafale lunaire sur trépied reste affectée par dérive et turbulence :
ne pas appeler « vérité » son empilement sans qualification.

Chronométrer chargement froid et calcul chaud séparément ; synchroniser le GPU
avant de lire le chronomètre, sinon on mesure parfois seulement une mise en file.
Inclure les transferts et le raccord des tuiles dans la latence utilisateur.
Annulation et résultat obsolète doivent empêcher toute publication partielle.

\chapter{Parcours pratique et exercices corrigés}

\section{Parcours en cinq ateliers}

Atelier 1 : exécuter la régression, changer un point et expliquer les gradients.
Atelier 2 : vérifier une dérivée par différences finies puis lire le code XOR.
Atelier 3 : tracer toutes les dimensions du U-Net et compter ses paramètres.
Atelier 4 : exécuter le TP PyTorch, changer une seule hypothèse et conserver
le témoin initial. Atelier 5 : préparer un dataset avec provenance et scènes
réservées avant de lancer une expérience photographique.

\section{Questions et solutions}

\paragraph{1. Combien de paramètres dans une couche dense $4\to3$ ?}
$4\times3+3=15$, avec un biais par sortie.

\paragraph{2. Et une convolution $3\times3$, $8\to16$ ?}
$9\times8\times16+16=1168$. Cela ne dépend pas de la largeur de l'image,
contrairement au coût de calcul et à la mémoire d'activations.

\paragraph{3. Pourquoi une sigmoïde finale gêne-t-elle le RGB signé ?}
Elle impose des sorties entre 0 et 1. Une observation négative ne peut plus
être représentée, même si l'architecture et les poids étaient parfaits.

\paragraph{4. Pourquoi force zéro ne veut-elle pas dire sigma zéro ?}
Un réseau à sigma nul peut encore modifier l'image. Une identité exacte
se garantit par le contrat de contournement, pas par un espoir sur les poids.

\paragraph{5. Peut-on partager les crops d'une photo entre train et test ?}
Pas pour mesurer une généralisation à de nouvelles scènes : textures,
éclairage et traitements restent communs. Garder la scène entière dans un groupe.

\paragraph{6. Pourquoi la variance d'une moyenne ne baisse-t-elle pas toujours en $1/N$ ?}
Les covariances ajoutent des termes :
$\operatorname{Var}(N^{-1}\sum X_i)=N^{-2}\sum_{ij}\operatorname{Cov}(X_i,X_j)$.
Le bruit commun ne disparaît pas avec une simple moyenne.

\paragraph{7. Pourquoi appeler \texttt{eval()} et \texttt{no\_grad()} ?}
Le premier règle le comportement de certains modules ; le second désactive
le suivi des gradients. Ils n'ont pas la même responsabilité.

\paragraph{8. Le TP apprend-il à débruiter nos RAW ?}
Non. Il ajuste une rampe artificielle avec bruit gaussien et vérifie le câblage.
Il n'a ni dataset Canon, ni référence photographique, ni validation indépendante.

\paragraph{9. Une baisse du bruit apparent suffit-elle ?}
Non : vérifier matière, fins contours, biais, couleur et hallucinations,
sinon un simple flou peut gagner le mauvais concours.

\paragraph{10. Pourquoi ne pas commencer par un modèle énorme ?}
Sans données et contrats fiables, sa capacité accélère aussi la mémorisation
des biais ; coût et complexité rendent le diagnostic plus difficile.

\chapter{Références et suite dans IPR}

Ce cours et son code sont une rédaction pédagogique originale. Les sources
ci-dessous servent à approfondir ; aucune promesse de reprendre leurs poids
ou leur implémentation n'est faite.

\begin{itemize}
\item PyTorch, \href{https://docs.pytorch.org/tutorials/beginner/basics/intro.html}{Learn the Basics} : modules, données et boucle d'apprentissage.
\item PyTorch, \href{https://docs.pytorch.org/tutorials/beginner/basics/autogradqs_tutorial.html}{Automatic Differentiation} : graphe et gradients automatiques.
\item Ronneberger, Fischer et Brox,2015, \href{https://lmb.informatik.uni-freiburg.de/people/ronneber/u-net/}{U-Net} : principe encodeur/décodeur et connexions latérales, initialement pour la segmentation.
\item Goodfellow, Bengio et Courville,2016, \href{https://www.deeplearningbook.org/}{Deep Learning} : approfondissement mathématique et méthodologique, ouvrage distinct de ce cours.
\item Brooks et al.,2019, \href{https://github.com/google-research/google-research/tree/master/unprocessing}{Unprocessing Images for Learned Raw Denoising} : importance des gains et du pipeline dans les données d'apprentissage.
\end{itemize}

La \href{../architecture/restauration-neuronale.md}{préparation B4} et ses JSON
décrivent séparément les observations matérielles, les besoins de corpus et
le protocole de qualification. Les14 tests du premier banc U-Net synthétique
ne constituent ni une validation des rendus, ni une livraison de la bêta.
Le choix d'un entraînement photographique doit attendre un contrat, des données
qualifiées et des critères de sortie explicites.

\end{document}
